\documentclass{article}

\usepackage[T1]{fontenc}

\usepackage{booktabs,graphicx,siunitx}

\begin{document}

\section{Multiple linear regression}

\begin{tabular}{llll}
\toprule
 & (1) & $\beta$ & VIF \\
\midrule
(Intercept) & 20.36 &  &  \\
 & (6.03) &  &  \\
 & [8.11, 32.62] &  &  \\
pre\_score & 0.61 & --- & 1.25 \\
 & (0.08) &  &  \\
 & [0.46, 0.77] &  &  \\
age & 0.02 & --- & 1.41 \\
 & (0.07) &  &  \\
 & [$-$0.12, 0.16] &  &  \\
group: b & 5.35 & --- & 1.04 \\
 & (1.64) &  &  \\
 & [2.01, 8.69] &  &  \\
method: online & $-$2.05 & --- & 1.25 \\
 & (2.02) &  &  \\
 & [$-$6.16, 2.06] &  &  \\
method: workshop & $-$3.25 & --- & 1.25 \\
 & (2.27) &  &  \\
 & [$-$7.87, 1.37] &  &  \\
\midrule
Num.Obs. & 40 &  &  \\
R\textsuperscript{2} & 0.75 &  &  \\
R\textsuperscript{2} Adj. & 0.71 &  &  \\
F & 20.42 &  &  \\
RMSE & 4.68 &  &  \\
AIC & 250.96 &  &  \\
BIC & 262.78 &  &  \\
Log.Lik. & $-$118.48 &  &  \\
\bottomrule
\end{tabular}

\includegraphics{figures/01_multiple-linear-regression/figure_residuals.png}

\includegraphics{figures/01_multiple-linear-regression/figure_coefficient-plot.png}

Each predictor's B is its effect holding the others constant; p and CI assess it. R\textsuperscript{2} is the joint variance explained; VIF flags predictors that overlap too much (multicollinearity). B is in each predictor's own units, so to compare which predictors matter most use the standardized coefficient ($\beta$), not B. Your significance threshold ($\alpha$) is 0.05.

The model explained R\textsuperscript{2}=.75 of the variance, F(5,34)=20.42, p $<$ .001; predictor pre\_score gave B=0.61, p $<$ .001.

Statistical basis: Ordinary least squares regression coefficients (B), SE, CI (R Core Team (2026). R: A Language and Environment for Statistical Computing. R Foundation for Statistical Computing, Vienna.); Standardized coefficients ($\beta$) (Lüdecke D, Ben-Shachar MS, Patil I, Makowski D (2020). "Extracting, Computing and Exploring the Parameters of Statistical Models using R." Journal of Open Source Software, 5(53), 2445.); Variance inflation factor (VIF) cutoff guidance (O’Brien, R. M. (2007). "A caution regarding rules of thumb for variance inflation factors." Quality \& Quantity, 41(5), 673-690.).

\end{document}